% master.tex : master-fil for projektet
% ------------------------------------------------------------------------------
% Dette er hovedfilen for projektet, hvori indhold fra alle input-filer (tekst,
% billeder, litteraturdatabaser, osv.) samles

% Dokumenttypen 'book' er valgt pga. dens mange fleksible indstillinger
% Se https://tex.stackexchange.com/a/36989/118167
\documentclass[12pt,a4paper,oneside,openright,danish]{book}

% Variabler, som bruges til automatisk at indsætte titel, forfattere, osv. på
% forsiden og titelbladet.
\def \projecttitle       {Eksamensopgaver-OPG1}
\def \projectsubtitle    {Computerstøttede beregninger}
\def \projecttheme       {CSB}
\def \projectdegree      {Matematik}
\def \projectperiod      {Forårssemesteret 2020}
\def \projectnumber      {P2}
\def \projectgroup       {}
\def \projectauthors     {
	Daniel Hasanovic			\\
	Jens Koldkur				\\
	Jonas Bach Christensen		\\
	Julie Havbo Lund			\\
	Mads Bjerregaard Kjær		\\
	Mathias Gammelgaard			\\
	Rasmus Jespersgaard
}
\def \projectsupervisors {
 
}

% Preamblet indeholder alle de indstillinger og makroer, som skal indsættes for
% hovedindholdet, og i denne skabelon samles det i filen aaumath.sty, som
% definerer en pakke, der kan indlæses med \usepackage.
\usepackage{aaumath}

% Dokumentets indhold indsættes mellem \begin- og \end-makroerne for
% 'document'-blokken
\begin{document}

% Dokumentets 'front matter' tælles ikke med ifm. antal sider og nummereres med
% romerske tal. Herunder hører f.eks. forsiden, titelbladet, forordet og
% indholdsfortegnelsen.
%\frontmatter
% incl/misc/frontpage.tex : rapportens forside
% ------------------------------------------------------------------------------


\backgroundsetup{
  scale = 1,
  angle=0,
  opacity=1,
  contents = {
    \includegraphics[width=\paperwidth,height=\paperheight]{fig/img/aau/waves.pdf}
  }
}
\BgThispage
\pdfbookmark[0]{Forside}{forside}
\begin{titlepage}
  \centering
  \phantom{}
  \vspace{2cm}

  % AAU-segl
  \begin{minipage}[c]{0.2\paperwidth}
    \centering
    \makebox[0pt]{
      % fig/tikz/aau-badge.tex : AAU-logo til forsiden
% ------------------------------------------------------------------------------

\begin{tikzpicture}
  % Tegn hvid cirkel og tilføj det gennemsigtige, blå logo ovenpå
  \node[circle,color=white,fill=white,minimum size=1.175\textwidth] at (0,0) {};
  \node at (0,0) {\includegraphics[width=\textwidth]{fig/img/aau/logo-circle.pdf}};
\end{tikzpicture}

    }
  \end{minipage}

  % Hovedindhold
  \vspace{4cm}
  {\fontfamily{bch}\selectfont
    \fboxsep0pt\colorbox{white}{
      \begin{minipage}{\textwidth}
        \centering
        \color{AAUblue1}

        \vspace{2em}
        {\Huge\bfseries\projecttitle}

        {\Large\bfseries\projectsubtitle}

        \bigskip
        \parbox{\textwidth}{\centering\large\projectauthors}

        \bigskip
        {\projectdegree}
        \vspace{2em}
      \end{minipage}
    }
  }

\end{titlepage}

%% incl/misc/contents.tex :
% ------------------------------------------------------------------------------

\pdfbookmark[0]{Indhold}{indhold}

% Indstillingerne i denne fil er grupperet, så de ikke påvirker andre dele
\begingroup

% Slå twoside fra midlertidigt for at undgår sideskift
\makeatletter
\@twosidefalse
\makeatother

% Placer indholdsfortegnelsen på sin egen side (bedst når den kun fylder en side)
\tableofcontents
\clearpage

% Placer oversigtslisterne på efterfølgende sider
%\let\clearpage\relax
%\listoffigures
%\listoftables
%\listofalgorithms
%\lstlistoflistings

\endgroup


% Dokumentets 'main matter' (hovedindhold) er der, hvor det meste indhold skal
% sættes ind. Sider og overskrifter er nummererede med arabiske tal.
\mainmatter

% Input-filer bør opdeles således, at hver fil svarer til et kapitel. Makroen
% \include indsætter et sideskift og indholdet fra den givne stil.

	\chapter{Det hængende kabel}
% God arbejdslyst xD
\begin{color}{AAUblue2}
%
Denne delopgave giver en grundlæggende forståelse for kvadraturregler.
Givet et kabel ophængt i to pæle af samme højde vil under påvirkning af tyngdekraften antage en form, der kan beskrives ved en kurve, som på dansk kaldes en kædelinje, der tager formen 
\begin{align*}
f(x)=h_0 + \lambda \text{cosh} \left( \frac{x}{\lambda} \right).
\end{align*}
Kablets massetæthed betegnes med $\rho$ og spændingen i kablet i dets laveste punkt med $T$, så $\lambda = \frac{T}{\rho}$. Afstanden som kablet synker i midten, betegnes $s$.
% 
\end{color}
\section*{1. Interpolerende kvadraturregler}
% 
%
\begin{color}{AAUblue2}
%
fisk
% 
\end{color}
\\\\
% 
%
%
%%%%%%%%%%%%%%%%%%%%%%%%%%%%%%%%
%%% Flot graf alla Julie     %%%
%%%%%%%%%%%%%%%%%%%%%%%%%%%%%%%%
%
\begin{figure}[h!]
\begin{center}
\begin{tikzpicture}[scale=10]
%
% Koordinater 
% ------------------------------------------------------
\coordinate (a) at (0.1,0.5);
\coordinate (b) at (0.7,0.4);
\coordinate (c) at (0.4,0.483);
%
% Funktion
% ------------------------------------------------------
\filldraw[thick,color=AAUred!15] (0.1,0.483) -- (0.7,0.483) -- (0.7,0.3) -- (0.1,0.3) -- (0.1,0.483);
\draw[very thick, color=AAUblue2] (a) .. controls (0.3,0.8) and (0.6,0.1) .. (b);  
%
%
\draw[thick](0.1,0.28) -- (0.1,0.32); % (a)
\draw[thick](0.7,0.28) -- (0.7,0.32); % (b)
\draw[thick](0.4,0.28) -- (0.4,0.32); % (m_1)
%
% Punkt 
% -------------------------------------------------------
\filldraw [black] (a) circle (0.1pt);  
\filldraw [black] (b) circle (0.1pt); 
\filldraw [black] (c) circle (0.1pt); 
%
\node at (0.1,0.22) (){$a$};
\node at (0.7,0.22) (){$b$};
\node at (0.4,0.22) (){$\frac{a+b}{2}$};
%
% 
%
% Koordinatsystemet 
% -------------------------------------------------------
\draw[thick,->] (-0.05,0.3) -- (0.85,0.3) node[anchor=south east]{$x$};
\draw[thick,->] (0,0.2) -- (0,0.7) node[anchor=north west]{$y$};
%
%
\end{tikzpicture}
  \caption{Illustration af Midtpunktsreglen.}
  \label{fig:midtpunkt}
\end{center}
\end{figure}



%
%%%%%%%%%%%%%%%%%%%%%%%%%%%%%%%%
%%% Flot graf alla Julie     %%%
%%%%%%%%%%%%%%%%%%%%%%%%%%%%%%%%
%
\begin{figure}[h!]
\begin{center}
\begin{tikzpicture}[scale=10]
%
% Koordinater 
% ------------------------------------------------------
\coordinate (a) at (0.1,0.5);
\coordinate (b) at (0.7,0.4);
\coordinate (c) at (0.4,0.483);
%
% Funktion
% ------------------------------------------------------
\filldraw[thick,color=AAUred!15] (a) -- (b) -- (0.7,0.3) -- (0.1,0.3) -- (a);
\draw[very thick, color=AAUblue2] (a) .. controls (0.3,0.8) and (0.6,0.1) .. (b);  
%
%
\draw[thick](0.1,0.28) -- (0.1,0.32); % (a)
\draw[thick](0.7,0.28) -- (0.7,0.32); % (b)
%\draw[thick](0.4,0.28) -- (0.4,0.32); % (m)
%
% Punkt 
% -------------------------------------------------------
\filldraw [black] (a) circle (0.1pt);  
\filldraw [black] (b) circle (0.1pt); 
%\filldraw [black] (c) circle (0.1pt); 
%
\node at (0.1,0.22) (){$a$};
\node at (0.7,0.22) (){$b$};
%\node at (0.4,0.22) (){$\frac{a+b}{2}$};
%
% 
%
% Koordinatsystemet 
% -------------------------------------------------------
\draw[thick,->] (-0.05,0.3) -- (0.85,0.3) node[anchor=south east]{$x$};
\draw[thick,->] (0,0.2) -- (0,0.7) node[anchor=north west]{$y$};
%
%
\end{tikzpicture}
  \caption{Illustration af Trapezreglen.}
  \label{fig:trapez}
\end{center}
\end{figure}


%
%%%%%%%%%%%%%%%%%%%%%%%%%%%%%%%%
%%% Flot graf alla Julie     %%%
%%%%%%%%%%%%%%%%%%%%%%%%%%%%%%%%
%
\begin{figure}[h!]
\begin{center}
\begin{tikzpicture}[scale=10]
%
% Koordinater 
% ------------------------------------------------------
\coordinate (a) at (0.1,0.5);
\coordinate (b) at (0.7,0.4);
\coordinate (c) at (0.4,0.483);
%
% Funktion
% ------------------------------------------------------
\filldraw[thick,color=AAUred!15] (a) -- (c) -- (b) -- (0.7,0.3) -- (0.1,0.3) -- (a);
\draw[very thick, color=AAUblue2] (a) .. controls (0.3,0.8) and (0.6,0.1) .. (b);  
%
%
\draw[thick](0.1,0.28) -- (0.1,0.32); % (a)
\draw[thick](0.7,0.28) -- (0.7,0.32); % (b)
%\draw[thick](0.4,0.28) -- (0.4,0.32); % (m)
%
% Punkt 
% -------------------------------------------------------
\filldraw [black] (a) circle (0.1pt);  
\filldraw [black] (b) circle (0.1pt); 
\filldraw [black] (c) circle (0.1pt); 
%
\node at (0.1,0.22) (){$a$};
\node at (0.7,0.22) (){$b$};
\node at (0.4,0.22) (){$\frac{a+b}{2}$};
%
% 
%
% Koordinatsystemet 
% -------------------------------------------------------
\draw[thick,->] (-0.05,0.3) -- (0.85,0.3) node[anchor=south east]{$x$};
\draw[thick,->] (0,0.2) -- (0,0.7) node[anchor=north west]{$y$};
%
%
\end{tikzpicture}
  \caption{Illustration af Simpsonsreglen.}
  \label{fig:simpsons}
\end{center}
\end{figure}
%
\section*{2. Formler}
% 
\begin{color}{AAUblue2}
%
Vis, at følgende gælder:
\begin{align*}
f \left( \frac{L}{2} \right) = h_0 + \lambda \text{cosh} \left( \frac{L}{2 \lambda} \right) = h, \\
f(0) = h_0 + \lambda = h - s.
\end{align*}
Vis, at disse formler giver ligningen 
\begin{align*}
\lambda \text{cosh} \left( \frac{L}{2 \lambda} \right) = \lambda + s.
\end{align*}
%
\end{color}
\\
% 
%
Da søjlerne er placeret i $x = \pm \frac{L}{2}$ og søjlerne har højden $h$ gælder følgende
\begin{align}
f\left(\frac{L}{2}\right)	= h_0 + \lambda \cosh \left(\frac{L}{2 \lambda}\right) = h
\end{align}
Søjlerne stå således symetrisk omkring $x = 0$ og da $s$ er afstanden som kæden flader, haves følgende
\begin{align}
f(0) =	 h_0 + \lambda 	\cosh (0)	= h_0 + \lambda = h-s
\end{align}
%
Ligning (1.1) og (1.2) kan give anledning til en ny ligning:
%
\begin{align}
h_0 + \lambda \cosh \left(\frac{L}{2 \lambda}\right) -\left(h_0 + \lambda \cosh \left(\frac{L}{2 \lambda}\right)\right) &= h_0 + \lambda -(h_0 + \lambda)\\
h_0 + \lambda \cosh \left(\frac{L}{2 \lambda}\right) - h &= h_0 + \lambda - (h - s)\\
\lambda \cosh \left(\frac{L}{2 \lambda}\right) - h &= \lambda - h + s\\
\lambda \cosh \left(\frac{L}{2 \lambda}\right) &= \lambda + s
\end{align}
\section*{3. Adaptiv kvadratturregler}
%
% 
%
\begin{color}{AAUblue2}
%
\textbf{Opgavebeskrivelse:} 
Forklar idéen i en sammensat adaptiv kvadraturregel.
% 
\end{color}
\\\\
% 
\textbf{Besvarelse:} 
Idéen i de adaptive kvadraturregler omhandler, hvordan en funktion kan behandles, når der ikke kan udvælges punkter til underinddeling på forhånd. Eksempelvis hvis funktionen ikke er kendt, men funktionsværdien kendes for en mængde af punkter.
Der findes to primære måder at lave disse underinddelinger: \\
1. Der vælges et antal $N$ underinddelinger af det oprindelige interval af formen $\left [  x_k, x_{k+1} \right ]$, hvor $a=x_0<x_1<x_2<\ldots<x_N$ herefter anvendes den sammensatte Simpsons regel beskrevet i opgave 3 hvor antallet af underintervaller i de oprindelige $N$ gennemgår kontinuerlig fordobling.
(af uransagelige årsager siger bogen at 5 eller 20 er de mest normale inddellinger, dette bliver ikke uddybet) \\
2. Den anden måde benytter ligeledes den sammensatte Simpsons regel her på intervallet, denne udføres igen på de nye delintervaller samt på begge halvdele af disse, såfremt man er inden for den ønskede nøjagtighed for et givet interval accepteres dette som delresultat og der arbejdes nu med næste delinterval.
Dette gentages indtil alle delintervaller er opdelt i et antal intervaller der bringer det inden for den ønskede fejl.
(side 75 bunden)
%

\section*{4.}
% 
\section*{(a) }

\section*{(b) Python script - Sekantmetoden}
%Skriv et Python script der implementerer sekantmetoden. Det kan gøres ved at modicere scriptet beregnet til Newtons metode. Sammenlig sekantmetoden med Newtons metode.

Forskellen mellem Sekantmetoden og Newtons metode er at fiiiiiiiiiiiiiiiiiiiiiiiiiiiiiiiiiiiiiiiiiiiiiiiiiiiiiiiiiiiiiiiiiiisk 

\section*{(c) Fixpunktligningen}
% Ligningen (6) kan omskrives til en xpunktligning. Gr rede for de to omskrivninger her


\section*{(d) }
	\chapter{Interpolation af funktioner}
% God arbejdslyst
% 
%%%%%%%%% INTRODUKTION %%%%%%%%%
\section*{Introduktion}
Dette dokument giver en grundlæggende forståelse for Lagrange interpolation samt numerisk integrering. 

Først defineres lagrange interpolation
%%%%%%%%%%% OPGAVE 1 %%%%%%%%%%%
\section*{Opgave 1}
\textbf{Opgavebeskrivelse:} Gennemgå teorien for Lagrange interpolation baseret på datapunkterne $x_0, \ldots, x_N$ og $f_0, f_1, \ldots , f_N$. \\
\textbf{Besvarelse:} 
Lagrange interpolation handler om at finde et polynomium $p(x_k)=f(x_k)$ og er således en måde at approksimere en funktionsværdi givet $x_0 \ldots x_N$ og $N+1 $  $y$ værdier.
Det ønskes således at finde et polynomium med interpolationsegenskaben $p_{x_j}=y_j$,$j=0,1,\ldots,N$.

Problemet omhandler således at finde et $N$'te grads polynomium $p$, der tager funktionsværdien for $f$ i en givet mængde punkter. Ideen kan således relateres til taylorserier og konstruktion af taylorpolynomier da disse ligeledes antager funktionsværdien i et givet udviklingspunkt.
Lagrange polynomier er defineret ved:
\begin{align*}
l_k(x)=\frac{\prod_{\substack{j=0 \\ {j \neq k}}}^{N}(x-x_j)}{\prod_{\substack{j=0 \\ {j \neq k}}}^{N}(x_k-x_j)}=\frac{(x-x_0) \cdots (x-x_{j-1}) \cdot (x-x_{j+1})\cdots(x-x_N)}{(x_j-x_0)\cdots (x_j-x_{j-1}) \cdot (x_j-x_{j+1})\cdots (x_j - x_N)}
\end{align*}
og et Lagrange-interpolerede andensgradspolynumium er således givet ved:
\begin{align*}
p(x)=f(x_0)\frac{(x-x_1)(x-x_2)}{(x_0-x_1)}+f(x_1)\frac{(x-x_0)(x-x_2)}{(x_1-x_0)(x_1-x_2)}+f(x_2)\frac{(x-x_0)(x-x_1)}{(x_2-x_0)(x_2-x_1)}
\end{align*}
Der eksisterer et entydigt Lagrange interpoleret polynomium for punkter givet ved $x_0, \ldots, x_N$ og $f_0, f_1, \ldots , f_N$.
\begin{proof}
Ved at finde sådanne polynomier, er deres eksistens bevist. Deres entydighed vil nu bevises. 
Antag, at $q_n$ og $p_n$ er to forskellige polynomier med grad $\leq n$, og at de begge interpolerer samme data. 
Så er graden af polynomiet $p_n - q_n$ også $\leq n$, og værdien af dette polynomie er $0$ ved $n+1$ datapunkter. 
Men en polynomie med grad $n$ har højst $n$ nuller, medmindre det er nulpolynomiet. 
Således er $p_n - q_n = 0$, hvormed $p_n = q_n$. 
\end{proof}
%%%%%%%%%%% OPGAVE 2 %%%%%%%%%%%    
\section*{Opgave 2}
\textbf{Opgavebeskrivelse:} Vi definerer funktionen $f$ ved
\begin{align*}
f(x)=x^2-\sin(10x), x \in \mathbb{R}
\end{align*}
Vi ser på intervallet $\left [0,3 \right ]$ og definerer punkterne $x_j = \frac{3j}{N}, j = 0, 1, \ldots, N$, så at vi har $N+1$ ækvidistante punkter. Vi definerer også $f_j=f(x_j), j=0,1,\ldots,N.$ Vi betegner med $p_N$ det interpolerende polynomium baseret på disse data.  
\\
\textbf{Besvarelse:} Dette er blot definitioner, vi skal benyttes til resten af delopgaven. 
%%%%%%%%%%% OPGAVE 3 %%%%%%%%%%%
\section*{Opgave 3}
\textbf{Opgavebeskrivelse:} Scriptet \textsc{Interpolation bin.py} giver en vurdering af den maksimale fejl 
\begin{align*}
\underset{x \in \left [0,3 \right ]}{\text{max}} \lvert f(x)-p_N(x) \rvert
\end{align*}
for forskellige værdier af $N$. Det beregner også en teoretisk vurdering af fejlen, se nedenfor. Gennemgå dette script i detaljer. \textbf{Bemærk} at dette script ikke bruger \textsc{numpy}, men er baseret på lister. Begrundelsen er at vi skal være i stand til at modificere dette script til at anvende \textsc{decimal modulet}. Efter gennemgangen skal I modificere scriptet til at generere plots for små værdier af $N$, der i samme figur viser grafen for $f(x)$, grafen for $p(x)$ og knudepunkterne.
\\
\textbf{Besvarelse:} I scriptet starter man med at definere Lagrange funktionen. Dette gøres ved hjælp af lister, og en for-løkke, hvor j ikke må være lig k. Den ganger resultaterne fra listen temp samme, og finder produktet.\\
Så findes Lagrange interpolation polynomium. Så opskrives parametrene og funktioenn der ønskes undersøgt. Programmet finder efterfølgende nogle x- og y-værdier samt afstanden mellem x-værdierne. Så udregner programmet approximationen og løsningen. Til sidst udregner programmet fejlen og printer noget data. 

(Obs. spørg ind til hvorfra "Error bound" delen opstår (linje 31-34), og hvordan man kommer frem til dette 

%%%%%%%%%%% OPGAVE 4 %%%%%%%%%%%
\section*{Opgave 4}
\textbf{Opgavebeskrivelse:} Lav en grafisk undersøgelse af fejlen $\lvert f(x)-p_N(x) \rvert$, for en række forskellige værdier af N baseret på en modificeret udgave af \textsc{Interpolation bin.py.}
\\
\textbf{Besvarelse:} Man kan se, at approximationen passer til funktionen, når N = 15,30 og 45. Hvis der indsættes andre værdier, eksempelvis 16, 31 og 46 kan der ses, at der opstår store udsving i yderpunkterne i definitionsmængden [0,3], hvilket også stemmer overens med forlæsningsnoterne fra kursusgang 5(s.4-5) hvor der kan ses, at Runge har bevist, at yderpunkterne divergerer
%%%%%%%%%%% OPGAVE 5 %%%%%%%%%%%
\section*{Opgave 5}
\textbf{Opgavebeskrivelse:} Brug Theorem 16 i afsnit 6.2 i [1] til at lave en teoretisk vurdering af fejlen
\begin{align*}
\underset{x \in \left [0,3 \right ]}{\text{max}} \lvert f(x)-p_N(x) \rvert.
\end{align*}
%% SE HJÆLP I OPGAVEFORMULERING %% JA TAK %%

%%%%% THEREOM 16 %&%%%%
\begin{align*}
f(x)-p(x) = \frac{(x-x_0)(x-x_1)\cdots(x-x_N)}{(N+1)!}f^{(N+1)}(\xi)
\end{align*}
\text{for ethvert} $\xi$ $\in$ $\left [a,b \right ]$
%%%%%%%%%%%%%%%%%%%%%%%
\\
%%%%% HINT 1/2 %&%%%%
Lad
\begin{align*}
L(x)=(x-x_0)(x-x_1)\cdots(x-x_N)
\end{align*}
%%%%%%%%%%%%%%%%%%%%%%%
\\
For ækvidistante punkter $x_{j+1} - x_j = h$ kan man vise, at der gælder
%%%%% HINT 2/2 %&%%%%
\begin{align*}
\lvert L(x) \rvert \leq \frac{1}{4}(N+1)!h^{N+1}, \text{   } x \in \left [x_0, x_N \right ] 
\end{align*}
%%%%%%%%%%%%%%%%%%%%%%%
\textbf{Besvarelse:} 
Udtrykket for maksimal fejl kan derfor omskrives til
\begin{align*}
\underset{x \in \left [0,3 \right ]}{\text{max}} \lvert f(x)-p_N(x) \rvert \leq \frac{{1}{4}(N+1)!h^{N+1}}{(N+1)!}f^{(N+1)}(\xi)=\frac{1}{4}h^{N+1}f^{(N+1)}(\xi)
\end{align*}
Det skal nu vises at denne er lig  $0.25h^{(N+1)}M$ for $M=10^{(N+1)}$ ved funktionen $f(x)=x^2-\sin(10x), x \in \mathbb{R}$ i intervallet $a=0$ til $b=3$ hvor $h$ altså antallet af steplings bliver $h=abs(b-a)/N$. Det skal nu vises at $f(x)^{(N+1)}(\xi)$ kan omskrives til $M=10^{(N+1)}$
for $N=0$ gælder: $f'(x)=2x-10cos(10x)$ for $n=1$ gælder at $f''(x)=2+100sin(10x)$ for $n=2$ gælder $f'''(x)=1000sin(10x)$ den afledede vil herefter skifte mellem sinus og cosinus og maksimalværdien for disse vil altid være 1 og den maksismale værdi for N'te afledte bliver derfor $10^{N+1}$.
\\
%%%%%%%%%%% OPGAVE 6 %%%%%%%%%%%
\section*{Opgave 6}
\textbf{Opgavebeskrivelse:} Hvor stor skal $N$ være, for at fejlen er mindre end $10^{-10}$? 
Brug den oprindelige version af \textsc{Interpolation bin.py} til at sammenligne teori og numerisk eksperiment. 
Hvad er konklusionen? 
Hvor kommer problemerne fra? 
I den forbindelse kan I inddrage den grafiske undersøgelse fra punkt 4.
\\
\textbf{Besvarelse:} Den teoretiske fejl (max) bliver $10^{-10}$ ved $N=48$, men i det numeriske eksperiment bliver fejlen aldrig mindre eller lig $10^{-10}$. Det problematiske er divergens i yderpunkterne, hvilket gør fejlen betydeligt større end den teoretiske, da de går mod $ \infty $ og $ - \infty $. Disse divergenspunkter kan også ses på grafen af $f(x)$ og $p(x)$. 
%%%%%%%%%%% OPGAVE 7 %%%%%%%%%%%
\section*{Opgave 7}
\textbf{Opgavebeskrivelse:} Gentag undersøgelserne af den maksimale fejl og sammenlign med den teoretiske vurdering ved at bruge decimalaritmetik-scriptet \textsc{Interpolation dec.py}. 
Prøv med forskellige værdier for antallet af betydende cifre. 
Hvad er konklusionen?
\\
\textbf{Besvarelse:} 
Ved lave værdier af betydende cifre er der dårlig præcision, og fejlen overstiger det teoretiske øvre estimat for grænsen. Vælges derimod et højt antal betydende cifre, eksempelvis $80$, bliver computationstiden højere, men fejlen forbliver mindre end grænsen for alle værdier af $N$.
Dette sker da der ved tidligere afrundinger skabes en situation hvori afrundringsfejlen bliver mere problematisk des højere værdi af $N$ idet der multipliceres med afrundingsfejl i flere iterationer.
%%%%%%%%%%% OPGAVE 8 %%%%%%%%%%%
\section*{Opgave 8}
\textbf{Opgavebeskrivelse:} Vi erstatter nu de ækvidistante punkter $x_j , j = 0, 1, \ldots , N$ med Chebyshevpunkterne $\Tilde{x}_j = 3 \frac{(1 -\cos(\frac{j \pi}{N}))}{2}, j = 0, 1, \ldots , N$. 
Gennemfør den grafiske undersøgelse fra punkt 4. ovenfor med dette valg af knudepunkter. 
Hvad kan man konkludere ud fra denne undersøgelse?
\\
\textbf{Besvarelse:} 
Ved valg af Chebyshevpunkterne som x-værdier sker der ikke de store udsving omkring grænseværdierne, og der ses tendens til at fejlen stablieres omkring $10^-{-14}$ ved $N=40$ og frem. Der er dog små udsving i værdien herefter og ikke et entydigt mønster i hvorvidt aproksimationen bliver bedre.
At aproksimationen stabliseres omkring yderpunkterne er endvidere i overenstemmelse med teorien, hvor disse netop er karakteriseret ved at minimere Runges fænomen.
%%%%%%%%% DONZO GONZO %%%%%%%%%%

% Appendicer indsættes inde i en appendices-blok og bliver nummereret med
% bogstaver i stedet for tal
\begin{appendices}
	\lstset{style=mystyle}
\lstinputlisting[language=Python]{code/delopgave1_1.py}
\end{appendices}

% Dokumentets 'back matter' er til ekstra ting som f.eks. litteraturlisten.
% Overskrifter bliver ikke nummereret her.
\backmatter

% Automatisk litteraturliste baseret på, hvilke kilder, der er blevet refereret
% til i løbet af rapporten.
\bibliographystyle{apalike}
\bibliography{
  incl/bib/books,
  incl/bib/articles,
  incl/bib/software
}

\end{document}
